\documentclass[../main.tex]{subfiles}
\begin{document}
\setlist[itemize]{topsep=2pt,itemsep=1pt,parsep=0pt,partopsep=0pt}
\setlist[itemize,2]{topsep=1pt,itemsep=0pt,parsep=0pt,partopsep=0pt}
\setlength{\parskip}{3pt}

\chapter{Week 3 — Research Material and Data}
\label{ch:week3}

\textit{Snapshot and part summaries: see the overview on page~\pageref{ov:week3}.}

\section{Part 1 — Population, Sample, and the Sampling Method}
\label{sec:sampling}

\begin{itemize}
  \item \keyterm{Population}: entire collection of things in the field — too large to study whole.
  \item \keyterm{Sample}: chosen subset; representativeness must be proven statistically with a required sample size. Non-representative samples limit conclusions to the studied portion only.
  \item Four-stage funnel: \keyterm{theoretical population} $\to$ \keyterm{study population} (accessible) $\to$ \keyterm{sampling frame} (listing mechanism) $\to$ \keyterm{sample} (units actually studied). Each stage can introduce hidden bias.
\end{itemize}

\section{Part 2 — Types of Samples and Standard Error}
\label{sec:sample-types}

\begin{itemize}
  \item \keyterm{Probability sampling}: every element has a known selection probability — supports statistical generalisation.
    \begin{itemize}
      \item \keyterm{Simple random}: equal chance for all possible samples.
      \item \keyterm{Complex probability}: systematic (every $n$th element), cluster (whole groups then within), or multi-stage.
    \end{itemize}
  \item \keyterm{Non-probability sampling}: no known probabilities — cannot claim representativeness and must state this explicitly.
    \begin{itemize}
      \item \keyterm{Convenience}: whatever is available.
      \item \keyterm{Purposive}: targeted for information richness (expert advice, quota, heterogeneity, snowball).
    \end{itemize}
  \item \keyterm{Standard error}: precision of the sample estimate, derived from the \keyterm{standard deviation} (spread around the mean). Larger spread $\uparrow$ error; larger sample $\downarrow$ error (\autoref{sec:eop}).
\end{itemize}

\section{Part 3 — Data Sources, Quality, and the Logbook}
\label{sec:data-sources}

\begin{itemize}
  \item Data sources: someone else's data, theoretical, simulated, or experimental (testing, surveying, observing).
  \item \keyterm{Personal data}: traceable to a natural person $\to$ GDPR applies; store in a secure EU database (DANS), not Google/Dropbox; delete or anonymise when no longer needed; report suspected leaks to the faculty data steward.
  \item Commercially sensitive data: follow company standards and check TU Delft open-access policy before signing anything.
  \item \keyterm{Protected data}: nuclear-weapons and missile-technology knowledge $\to$ knowledge-embargo exemption required (UN/EU).
  \item \keyterm{Qualitative data}: text, photos, environmental observations. \keyterm{Quantitative data}: numerical measurements or simulations. Type follows from strategy and variables; conversion between types is possible.
  \item \keyterm{Data quality}: check for typing errors, outliers, and coding errors before analysis.
  \item \keyterm{Research logbook}: records design choices, assumptions, experiment procedures, incidents, file locations and names, and people consulted — the traceability instrument formalised in Week~5 (\autoref{sec:dmp}).
\end{itemize}

\section{Reading Material}
\label{sec:reading3}

\begin{partbox}
The Week~3 reading is the TU Delft Data Usage and Storage Rules, specifying procedures, contacts, and storage requirements for the three data categories named in the slides.
\end{partbox}

\begin{itemize}
  \item Personal data: follow HREC guidelines, collect minimally, use special-category care (race, religion, health), keep within EU, delete/anonymise when done.
  \item Commercially sensitive: comply with company standards and TU Delft open-access policy before signing.
  \item Protected fields (hypersonic aerodynamics, re-entry, GNC, launch structures): knowledge-embargo exemption can take up to three months — plan accordingly.
\end{itemize}

\end{document}
